La planeación es la primera función del proceso administrativo porque permite definir hacia dónde quiere avanzar una organización y qué acciones son necesarias para alcanzar sus metas. En términos simples, consiste en establecer objetivos, anticipar necesidades, asignar recursos y fijar tiempos para orientar el trabajo futuro. Diversas obras de administración coinciden en que la planeación es una herramienta esencial para convertir una intención general en decisiones concretas y coordinadas \parencite{robbins2017,bateman2020}.

Desde una perspectiva empresarial, la planeación adquiere una importancia mayor porque una empresa manufacturera o de servicios debe coordinar producción, inventarios, personal, costos, tiempos y calidad. Si no se planea adecuadamente, aparecen pérdidas por retrasos, desperdicio de materiales, mala distribución de tareas y decisiones impulsivas. Por eso, la planeación no solo organiza el trabajo, sino que también reduce la incertidumbre, mejora la toma de decisiones y fortalece la capacidad de responder a cambios del mercado, la competencia o la demanda.

En un entorno competitivo, la planeación permite que la empresa se mantenga orientada hacia sus objetivos y evite la improvisación. La organización conoce qué debe hacer, cómo lo hará, cuándo y con qué recursos. Esto genera claridad, orden y mayor probabilidad de éxito, especialmente en actividades complejas donde varias áreas deben actuar de manera sincronizada. En ciudades donde muchas empresas dependen de la industria, la logística y el comercio, una buena planeación ayuda a sostener la operación diaria y a responder con mayor rapidez a los cambios del entorno.

Como lo resalta la experiencia práctica, no hay viento favorable para un barco sin rumbo. Por ello, es necesario saber hacia dónde debe ir la empresa u organización, y este propósito se logra precisamente con una buena planeación.